% Fluxo em blocos da estrutura laboratorial do DEE.
% Layout 2x2 com eixo central: o distribuidor "MACROÁREAS" desce por um corredor
% livre entre as colunas e alimenta as quatro caixas por derivações curtas,
% sem contornar o desenho e sem sobreposição entre caixas.

\begin{center}
\resizebox{0.96\textwidth}{!}{%
\begin{tikzpicture}[
  >=stealth,
  topo/.style={
    draw=azulpetroleo,
    fill=azulpetroleo!9,
    very thick,
    rounded corners=2mm,
    align=center,
    text width=10.4cm,
    minimum height=1.20cm,
    inner sep=3mm,
    font=\sffamily\bfseries
  },
  centro/.style={
    draw=azulpetroleo!70,
    fill=white,
    thick,
    rounded corners=1.8mm,
    align=center,
    text width=3.5cm,
    minimum height=0.85cm,
    inner sep=2.5mm,
    font=\sffamily\bfseries\small
  },
  area/.style={
    draw=azulpetroleo!80,
    fill=cinzaclaro,
    thick,
    rounded corners=2mm,
    align=left,
    text width=7.0cm,
    minimum height=4.0cm,
    inner sep=4mm,
    font=\small
  },
  uso/.style={
    draw=ifmaverde!75!black,
    fill=ifmaverde!6,
    thick,
    rounded corners=2mm,
    align=center,
    text width=15.8cm,
    minimum height=1.35cm,
    inner sep=3.5mm,
    font=\footnotesize
  },
  eixo/.style={thick,color=azulpetroleo!75},
  seta/.style={->,thick,color=azulpetroleo!75},
  no/.style={circle,fill=azulpetroleo!75,inner sep=0pt,minimum size=1.6mm}
]

% Bloco principal
\node[topo] (dee) at (0,0) {Departamento de Eletroeletrônica -- DEE\\[-0.2mm]
\normalfont\footnotesize IFMA -- Campus São Luís/Monte Castelo};

% Nó de distribuição
\node[centro] (macro) at (0,-1.8) {MACROÁREAS};
\draw[seta] (dee.south) -- (macro.north);

% Primeira linha (centro em y = -5,0; corredor central livre de 1,4 cm)
\node[area] (auto) at (-4.6,-5.0) {
\centering\sffamily\bfseries AUTOMAÇÃO\par\vspace{2mm}
\raggedright\normalfont
$\bullet$ Automação Industrial\\[0.8mm]
$\bullet$ Redes Industriais e IoT\\[0.8mm]
$\bullet$ Robótica\\[0.8mm]
$\bullet$ Pneumática e Eletropneumática\\[0.8mm]
$\bullet$ Instrumentação e Sensoriamento};

\node[area] (elet) at (4.6,-5.0) {
\centering\sffamily\bfseries ELETRÔNICA E TELECOMUNICAÇÕES\par\vspace{2mm}
\raggedright\normalfont
$\bullet$ Eletrônica Analógica e Digital\\[0.8mm]
$\bullet$ Sistemas Digitais e Embarcados\\[0.8mm]
$\bullet$ Telecomunicações, RF e Redes\\[0.8mm]
$\bullet$ Instrumentação Eletrônica\\[0.8mm]
$\bullet$ Eletrônica de Potência};

% Segunda linha (centro em y = -10,2; 1,2 cm de folga para a primeira linha)
\node[area] (eletrot) at (-4.6,-10.2) {
\centering\sffamily\bfseries ELETROTÉCNICA\par\vspace{2mm}
\raggedright\normalfont
$\bullet$ Instalações Elétricas\\[0.8mm]
$\bullet$ Medidas Elétricas\\[0.8mm]
$\bullet$ Máquinas Elétricas\\[0.8mm]
$\bullet$ Comandos Elétricos\\[0.8mm]
$\bullet$ Acionamentos e Qualidade de Energia};

\node[area] (sep) at (4.6,-10.2) {
\centering\sffamily\bfseries SISTEMAS DE ENERGIA\par\vspace{2mm}
\raggedright\normalfont
$\bullet$ Proteção e Subestações\\[0.8mm]
$\bullet$ Sistemas Elétricos de Potência\\[0.8mm]
$\bullet$ Smart Grid\\[0.8mm]
$\bullet$ Geração Distribuída\\[0.8mm]
$\bullet$ Análise e Monitoramento de Sistemas};

% Uso compartilhado
\node[uso] (uso) at (0,-14.0) {\textbf{Uso compartilhado da infraestrutura:} ensino técnico e superior $\bullet$ aulas práticas $\bullet$ iniciação científica $\bullet$ TCCs e monografias $\bullet$ pesquisas de mestrado e doutorado $\bullet$ extensão $\bullet$ capacitação profissional $\bullet$ projetos de P\&D, prototipagem, ensaios e validação tecnológica.};

% Eixo central: desce pelo corredor entre as colunas e termina no uso compartilhado.
\draw[eixo] (macro.south) -- (0,-12.9);
\draw[->,thick,color=ifmaverde!65!black] (0,-12.9) -- (uso.north);

% Derivações curtas para cada macroárea.
\foreach \y/\esq/\dir in {-5.0/auto/elet,-10.2/eletrot/sep}{%
  \node[no] at (0,\y) {};
  \draw[seta] (0,\y) -- (\esq.east);
  \draw[seta] (0,\y) -- (\dir.west);
}

\end{tikzpicture}%
}
\end{center}
